\section{Conclusion}

In this work, we propose a river valley landscape metaphor to explain why large learning rates can make implicit progress in training that can only be revealed when the learning rate is decayed. Further, we propose a simple data model that can generate the river valley landscape, attributing the difference of sharpness in different directions to the difference in uncertainty between different tokens. Based on our theory, we propose a learning rate schedule called $\wsds$ that can produce multiple intermediate checkpoints in a single-consecutive run as good as if we prespecified the individual budgets and used a tuned cosine learning rate. The learning rate schedule involves rapidly decaying to a minimal learning rate level at these budgets and resuming to a high constant learning rate immediately afterward and keeping the learning rate constant in the rest of the runs. From training language models from 0.1B to 1.2B parameters, we show $\wsds$ performs competitively with cosine learning rate and outperforms other compute-agnostic learning rates including $\wsd$ and cyclic-cosine learning rate fixing the total compute budget.

\section*{Acknowledgement}

The authors would like to thank the support of NSF 2211780 and also would like to thank the Google TPU Research Cloud for the computing resources that enabled these experiments